\textbf{目的：}围绕甲状腺乳头状微小癌（papillary thyroid microcarcinoma, PTMC）患者术前中央区淋巴结转移（CentralLNM）风险评估这一核心临床问题，在统一患者级 5-fold 交叉验证协议下，重新梳理原始临床基线模型、传统影像组学方法、多模态深度学习融合模型及其扩展验证结果，形成符合临床期刊逻辑的主文框架。\\
\textbf{方法：}首先基于原始术前临床变量构建患者级 Logistic 回归基线；随后在同一患者级划分下比较传统影像组学两条路径（ROI Radiomics-only、Clinical+ROI Radiomics）及其多模型表现；在深度学习部分，采用多模态超声输入、患者级 Attention-MIL、临床主导门控融合与先验约束损失构建主模型，并将主结果统一替换为 ViT-B/16 对应版本；进一步开展骨干网络公平比较、关键模块消融和三源融合验证。主要终点为 CentralLNM 的患者级测试集 AUC 均值$\pm$标准差，同时报告 ACC、F1、Sensitivity、Specificity；对于主模型额外报告逻辑依从率与风险分层表现。\\
\textbf{结果：}原始临床基线在 CentralLNM 上达到 AUC 0.6859$\pm$0.0658。ROI radiomics 中，ROI Radiomics-only 表现较弱（AUC 0.5756$\pm$0.0538），Clinical+ROI Radiomics 取得 AUC 0.6923$\pm$0.0696，提示组学相对原始临床仅带来有限增量。多模态深度学习主模型采用 ViT-B/16 后，在 CentralLNM 上达到当前最高 AUC 0.7104$\pm$0.0253（ACC 0.6289，F1 0.6686，Sen 0.6969，Spe 0.5490），较原始临床提高 0.0245，且折间波动明显下降。在剔除 Base-A（Clinical-only）后，消融结果仍显示 Ours 方案在 AUC 与逻辑一致性上保持优势。三源融合验证中，ViT-B/16 的 V1（Deep+Clinical）优于 V2（Tri-Fusion）（AUC 0.7104$\pm$0.0253 vs 0.6404$\pm$0.0496），提示当前三源注入策略仍需优化。\\
\textbf{结论：}在当前数据与评估协议下，原始临床信息仍是最稳定的预测基础；传统影像组学与深度影像分支均能提供补充信息，但其增益受融合方式与模型稳定性影响。患者级 MIL、临床主导门控与先验约束的联合设计，能够在保持临床逻辑一致性的同时提升深度模型的稳定性，可作为后续外部验证与临床转化研究的主要候选框架。
